% Reader question: which comparison detects a delegated permission restored
% after its parent removed it?
% Claim: root-only checking accepts the same chain that per-parent checking
% rejects. Authentication of every link remains a separate assumption.
% Evidence: analyses/harbor_card_v6_multihop_attack.pv and
% analyses/harbor_card_v7_multihop_fixed.pv, committed ProVerif results.
% Grammar: two aligned chain panels, unchanged sets, different comparison edges.
\begin{figure}[H]
\centering
\pdfigurehue{pdcobalt}%
\begin{tikzpicture}[pd figure,x=1cm,y=1cm]
  \node[pd title,anchor=west] at (0,5.15) {A. Compare the leaf with the root};
  \node[pd state,minimum width=2.4cm] (A1) at (1.3,4.25) {$A:\{r,w\}$};
  \node[pd state,minimum width=2.4cm] (B1) at (5.0,4.25) {$B:\{r\}$};
  \node[pd state,minimum width=2.4cm] (C1) at (8.7,4.25) {$C:\{r,w\}$};
  \draw[pd arrow] (A1) -- (B1);
  \draw[pd arrow] (B1) -- (C1);
  \draw[pd breach rule] (A1.south) -- (1.3,3.3) -- (8.7,3.3) -- (C1.south);
  \node[pd tag] at (5.0,3.3) {$C\subseteq A$: passes};
  \node[pd label,anchor=north] at (5.0,2.9)
    {accepts a write permission that $B$ did not hold};

  \node[pd title,anchor=west] at (0,1.95) {B. Compare every child with its parent};
  \node[pd state,minimum width=2.4cm] (A2) at (1.3,1.05) {$A:\{r,w\}$};
  \node[pd state,minimum width=2.4cm] (B2) at (5.0,1.05) {$B:\{r\}$};
  \node[pd state,minimum width=2.4cm] (C2) at (8.7,1.05) {$C:\{r,w\}$};
  \draw[pd arrow] (A2) -- (B2);
  \draw[pd arrow] (B2) -- (C2);
  \draw[pd focus rule] (A2.south) -- (1.3,.1) -- (4.85,.1) -- (4.85,.68);
  \node[pd tag] at (3.1,.1) {$B\subseteq A$: passes};
  \draw[pd breach rule] (5.15,.68) -- (5.15,.1) -- (8.7,.1) -- (C2.south);
  \node[pd tag] at (7.0,.1) {$C\nsubseteq B$: fails};
  \node[pd label,anchor=north] at (5.0,-.35)
    {rejects the attempted restoration of write permission};
\end{tikzpicture}
\caption{The capability sets are the same in both panels; only the checks
change. If every link is authenticated, checking each child against its
immediate parent stops re-escalation. ProVerif checks this two-link chain;
subset transitivity covers longer chains.}
\label{fig:anchor-v6v7-escalation}
\end{figure}
